% Options for packages loaded elsewhere
\PassOptionsToPackage{unicode}{hyperref}
\PassOptionsToPackage{hyphens}{url}
%
\documentclass[
]{article}
\usepackage{amsmath,amssymb}
\usepackage{iftex}
\ifPDFTeX
  \usepackage[T1]{fontenc}
  \usepackage[utf8]{inputenc}
  \usepackage{textcomp} % provide euro and other symbols
\else % if luatex or xetex
  \usepackage{unicode-math} % this also loads fontspec
  \defaultfontfeatures{Scale=MatchLowercase}
  \defaultfontfeatures[\rmfamily]{Ligatures=TeX,Scale=1}
\fi
\usepackage{lmodern}
\ifPDFTeX\else
  % xetex/luatex font selection
\fi
% Use upquote if available, for straight quotes in verbatim environments
\IfFileExists{upquote.sty}{\usepackage{upquote}}{}
\IfFileExists{microtype.sty}{% use microtype if available
  \usepackage[]{microtype}
  \UseMicrotypeSet[protrusion]{basicmath} % disable protrusion for tt fonts
}{}
\makeatletter
\@ifundefined{KOMAClassName}{% if non-KOMA class
  \IfFileExists{parskip.sty}{%
    \usepackage{parskip}
  }{% else
    \setlength{\parindent}{0pt}
    \setlength{\parskip}{6pt plus 2pt minus 1pt}}
}{% if KOMA class
  \KOMAoptions{parskip=half}}
\makeatother
\usepackage{xcolor}
\usepackage[margin=1in]{geometry}
\usepackage{graphicx}
\makeatletter
\def\maxwidth{\ifdim\Gin@nat@width>\linewidth\linewidth\else\Gin@nat@width\fi}
\def\maxheight{\ifdim\Gin@nat@height>\textheight\textheight\else\Gin@nat@height\fi}
\makeatother
% Scale images if necessary, so that they will not overflow the page
% margins by default, and it is still possible to overwrite the defaults
% using explicit options in \includegraphics[width, height, ...]{}
\setkeys{Gin}{width=\maxwidth,height=\maxheight,keepaspectratio}
% Set default figure placement to htbp
\makeatletter
\def\fps@figure{htbp}
\makeatother
\setlength{\emergencystretch}{3em} % prevent overfull lines
\providecommand{\tightlist}{%
  \setlength{\itemsep}{0pt}\setlength{\parskip}{0pt}}
\setcounter{secnumdepth}{-\maxdimen} % remove section numbering
\ifLuaTeX
  \usepackage{selnolig}  % disable illegal ligatures
\fi
\IfFileExists{bookmark.sty}{\usepackage{bookmark}}{\usepackage{hyperref}}
\IfFileExists{xurl.sty}{\usepackage{xurl}}{} % add URL line breaks if available
\urlstyle{same}
\hypersetup{
  pdftitle={estimating\_flux},
  pdfauthor={Jeff Wesner},
  hidelinks,
  pdfcreator={LaTeX via pandoc}}

\title{estimating\_flux}
\author{Jeff Wesner}
\date{2024-04-23}

\begin{document}
\maketitle

\hypertarget{purpose}{%
\subsection{Purpose}\label{purpose}}

The goal of this document is to demonstrate the workflow for estimating
global fluxes of contaminants. The are five primary data requirements.

\begin{enumerate}
\def\labelenumi{\arabic{enumi}.}
\tightlist
\item
  The global amount of annual insect emergence (mgDM/y)
\item
  The proportion of emergence that is carrying a given contaminant
\item
  The effect of the contaminant on insect emergence
\item
  The concentration of the contaminant in the surviving adult aquatic
  insects
\item
  The global amount of contaminant flux from adult aquatic insects
\end{enumerate}

A sixth data need that is not addressed here is the spatial location of
this flux (i.e., \emph{where} does the flux happen on earth?). This will
come from the spatial mapping approaches.

\hypertarget{approach}{%
\subsection{Approach}\label{approach}}

\begin{figure}
\includegraphics[width=1\linewidth]{plots/schematic} \caption{Figure 1. Schematic approach of the steps needed to estimate global contamiant (or nutrient) flux by adult aquatic insect emergence.}\label{fig:unnamed-chunk-1}
\end{figure}

Our overall approach (Figure 1) is to use literature estimates and
Bayesian modeling to quantify the uncertainty in each of the data needs
above, and then combine that uncertainty into estimates of the global
flux of contaminants transported by emerging aquatic insects from rivers
each year. This approach focuses on contaminants, but nutrient flux can
be added using the same methods.

We address 1-5 above (Figure 1) as follows, using mercury as a case
study:

\hypertarget{the-global-amount-of-annual-insect-emergence-mgdmy}{%
\subsubsection{1. The global amount of annual insect emergence
(mgDM/y)}\label{the-global-amount-of-annual-insect-emergence-mgdmy}}

The original data set contained X records of total community
macroinvertebrate annual production. A smaller number included
information on only insect community production, which is the measure we
are interested in, since it is only the insects that emerge. To obtain
estimates of insect production from the community production data, we
first estimated the proportion of community production that was insects
using an intercept only Beta model with the raw proportions as the
response variable (using only data that had both total community and
insect-only production). We then multiplied the posterior of the Beta
model by estimates of total community production, yielding modeled
estimates of aquatic insect secondary production (mean and sd) for each
collection.

Next, we estimated the fraction of this annual insect production that
emerged each year using data published in Gratton et al.~(2008). We fit
a random effects model to the Gratton data with varying intercepts for
taxon and study.

We then used a Monte Carlo approach to multiple the estimates of annual
insect secondary production by the fraction emerging each year
(\textasciitilde0.19), yielding n = 173 mean +/- sd estimates of insect
emergence (mgDM/m2/y). Based on the location of those estimates, we
assigned values for mean annual air temperature, elevation, and other
predictor variables from the BIOCLIM data set.

To estimate emergence as a function of temperature, we fit a Bayesian
generalized linear mixed model with emergence as the response variable
(posterior mean weighted by it's sd), temperature as the predictor, and
study (author/year) as a varying intercept. The likelihood was Gamma,
since the data were positive and continuous. This model revealed a weak
positive relationship between emergence and mean annual air temperature.

We then used the emergence vs temperature model to estimate total
emergence in the \(\sim 47000\) river basins reported by Allen and
Pavelsky (2018). We assigned a mean annual air temperature to each basin
and then used the posterior predictive distribution of the Gamma
regression to predict annual emergence in a based in units of mgDM/m2/y.
Then, we mutiplied that prediction by the area of river water in each
based to yield a total basin flux of mgDM/y. Summing that across basins
yields a global estimate of insect emergence from rivers each year:

\begin{verbatim}
## # A tibble: 5 x 2
##   statistic    `MT/dm/y`
##   <chr>            <dbl>
## 1 median_MTdmy   472631.
## 2 mean_MTdmy     540148.
## 3 sd_MTdmy       298869.
## 4 low95_MTdmy    182136.
## 5 high95_MTdmy  1277135.
\end{verbatim}

The results indicate that rivers export a mean of \(\sim 540000\) metric
tons (MT) of emerging insect dry mass each year with a 95\% credible
interval of 182,000 to 1,270,000 MTdm/y.

\hypertarget{the-proportion-of-emergence-that-is-carrying-a-given-contaminant}{%
\subsubsection{2. The proportion of emergence that is carrying a given
contaminant}\label{the-proportion-of-emergence-that-is-carrying-a-given-contaminant}}

Of the \(\sim 540000\) MTdm/y emergence, some fraction is carrying
contaminants. For example, if 2\% of surface river water is contaminated
with Hg, then we might assume that 2\% of emergence, or
0.02\$\times\$540,000 = 10,800 MTdm/y will contain Hg. We are not aware
of direct estimates of contaminant coverage globally.

\emph{Example for Mercury} To estimate the proportion of river area (and
hence insect emergence) impacted by Hg, we used Peterson et al.~Environ.
Sci. Technol. 2007, 41, 1, 58--65. Using random surveys of fish from
western US rivers, they reported that ``11\% of the assessed stream
length'' contained fish with more than 0.1 \(\mu\)gHg/g. Anything below
this level is consided save for consumption by fish-eating mammals.
Thus, we made the simplifying assumption that 11\% of stream lengths
(and hence area) would also result in Hg bioaccumulation in insects.
However, Peterson et al.~(2007) measured only streams \(\leq\) 60m wide.
Therefore, we needed to also assess how much of global stream widths are
\$\leq\$60m wide so that we can model Hg in a similar set of streams
(NOTE: Perhaps we can assume that Hg is also accumulating in larger
rivers, but here we assume that anything greater than 60m has zero Hg).

To assess the proportion of stream widths globally that are
\textgreater60m wide, we used the Pareto distrbution equation from Allen
and Pavelsky (2018). The estimated that stream widths globally are
distrbuted as a power law with a Pareto shape parameter of \$\sim\$0.83.
We simulated 10,000 stream widths using the \texttt{rpareto()} function
from the \texttt{VGAM} package, and calculated the proportion of stream
widths \$\leq\$60m. The result indicates that 99\% of stream widths are
likely less than 60m wide.

Hence, we assumed that the global flux of insects carrying Hg was/
\(540000 x 0.99 x 0.11 = 58806\)MTdm/y.

\hypertarget{the-effect-of-the-contaminant-on-insect-emergence}{%
\subsubsection{3. The effect of the contaminant on insect
emergence}\label{the-effect-of-the-contaminant-on-insect-emergence}}

For Hg, we did not assume any effect on insect emergence. For
contaminants that do have an effect on insect emergence, we would add an
additional multiplier to equation above (with uncertainty included as
well). For example, if a contaminant is expected to reduce emergenc by
0.5, then the value above would \(58806 x 0.5 = 29503\)MTdm/y

\hypertarget{the-concentration-of-the-contaminant-in-the-surviving-adult-aquatic-insects}{%
\subsubsection{4. The concentration of the contaminant in the surviving
adult aquatic
insects}\label{the-concentration-of-the-contaminant-in-the-surviving-adult-aquatic-insects}}

To estimate the concentration of contaminants in adult aquatic insects,
we used literature extractions for X contamintants that reported either
adult aquatic insect tissue concentrations, water concentrations, and/or
sediment concentrations.

\emph{Example for mercury} Using the 52(ish) measures of Hg in adult
aquatic insects, we ran a generalized linear model with adult Hg
concentration as the response variable, sediment Hg concentration as the
predictor variable, and study as a varying intercept. It suggested that
the mean Hg concentration in adult aquatic insects is \textasciitilde0.2
HgngDM/g

\hypertarget{the-global-amount-of-contaminant-flux-from-adult-aquatic-insects}{%
\subsubsection{5. The global amount of contaminant flux from adult
aquatic
insects}\label{the-global-amount-of-contaminant-flux-from-adult-aquatic-insects}}

Given the posterior predictions of the amount of potentially
contaminated insect emergence (\textasciitilde58,806 MTdm/y) along with
the posterior predictions of Hg concentration in those insect, we
estimate that the mean annual flux of Hg in emerging aquatic insects is:

\begin{verbatim}
## # A tibble: 5 x 7
##   proportion_with_mercury global_flux_mercury_kg_y .lower .upper .width .point
##                     <dbl>                    <dbl>  <dbl>  <dbl>  <dbl> <chr> 
## 1                    0.01                      1.3    0.4    4.2   0.95 median
## 2                    0.05                      6.5    2.2   21.2   0.95 median
## 3                    0.11                     14.4    4.8   46.6   0.95 median
## 4                    0.17                     22.2    7.4   72     0.95 median
## 5                    0.23                     30.1    9.9   97.4   0.95 median
## # i 1 more variable: .interval <chr>
\end{verbatim}

The table above gives five scenarios for the proportion of global
emergence impacted by Hg, from 1\% to 23\%. This provides a range of
scenarios centered around Peterson et al.'s estimate of 11\%. The
results suggest that emerging aquatic insects transport between 1 and 30
kg of Hg per year. For perspective, Pacific salmon of North America
transport \textasciitilde{} 7kg Hg/yr from marine to freshwater (Brandt
et al.~in review). Hence, these estimates appear reasonable to at least
an order of magnitude.

\end{document}
